% ==============================================================
%  Chapter 2 -- Theoretical Framework
% ==============================================================

\chapter{Theoretical Framework}
\label{ch:framework}

This chapter introduces the theoretical background on which the rest of
the thesis relies. We start with Decision Support Systems and their graphical formulation as Influence Diagrams, fix the formal
definitions that are used in subsequent chapters and motivate the
parameter-estimation problem that we tackle. We then present
Estimation of Distribution Algorithms, justify our design choices, including the specific
variants and the convergence guarantees that ground the empirical study of Chapter~\ref{ch:experiments}.

% ==============================================================
\section{Decision Support Systems and Influence Diagrams}
% ==============================================================

A \emph{Decision Support System} (DSS) is a model or piece of
software designed to assist decision-makers by facilitating access to
relevant data and by applying analytical techniques to
poorly structured problems~\cite{insua2005,shenoy2009eolss}.
The role of a DSS is not to replace the decision-maker but to provide
structured access to information and a systematic way of evaluating
alternatives before a final decision is reached.

Among the many graphical formalisms available, \emph{Influence
Diagrams} (IDs) have established themselves as one of the most
compact and expressive
representations~\cite{howard2005influence,koller2009pgm}. An ID
is a directed acyclic graph (DAG) whose nodes encode the three fundamental
ingredients of a sequential decision problem: random variables,
controllable choices and preferences, and whose arcs encode
conditional dependencies and informational precedence.

\subsection{Formal definition}

\begin{definition}[Influence Diagram]
\label{def:id}
An \emph{Influence Diagram} is a tuple
$\mathcal{D} = (\mathbf{C}, \mathbf{D}, V, \mathcal{E})$ where:
\begin{itemize}
    \item $\mathbf{C} = \{C_1,\ldots,C_n\}$ is a finite set of
          \emph{chance nodes}, each modelling a discrete random
          variable with finite domain $\Omega_{C_i}$;
    \item $\mathbf{D} = \{D_1,\ldots,D_k\}$ is a finite set of
          \emph{decision nodes}, each with a finite action set
          $\mathcal{A}_{D_i}$;
    \item $V$ is a single \emph{utility} (value) node encoding the
          preferences of the decision-maker;
    \item $\mathcal{E}$ is a set of directed arcs split into
          probabilistic, informational and functional arcs.
\end{itemize}
\end{definition}

For the ID to admit a well-defined evaluation procedure the graph
must be \emph{regular}. That is, a DAG where the value node has no successors and the decisions can be linearly ordered by a directed path that traverses all of them. Throughout this work all IDs are
assumed to be regular.

\subsubsection*{A concrete example: the bypass influence diagram}
\label{sec:bypass-id-intro}

We illustrate every notion introduced below on the \emph{bypass}
influence diagram, a small clinical model for the management of
coronary heart disease that we use throughout the rest of the
thesis. It contains six chance nodes
(\texttt{HeartDisease}, \texttt{Angiogram}, \texttt{Pain},
\texttt{EarlyResults}, \texttt{LifeQ}, \texttt{Economicalc}), two
decision nodes (\texttt{HeartSurgery}, \texttt{HeartPharma}) and a
single utility node. Its qualitative structure is shown in
Figure~\ref{fig:bypass-id}.

\begin{figure}[h]
\centering
\begin{tikzpicture}[
    every node/.style={font=\sffamily\scriptsize, align=center},
    chance/.style={ellipse, draw=blue!50, fill=cyan!10,
                   minimum width=2.4cm, minimum height=0.9cm},
    decision/.style={rectangle, draw=black, fill=orange!30,
                     minimum width=2.6cm, minimum height=0.9cm},
    utility/.style={regular polygon, regular polygon sides=6,
                    shape border rotate=90,
                    draw=black, fill=purple!30,
                    minimum width=2.6cm, minimum height=1.1cm,
                    inner sep=2pt},
    arr/.style={-{Stealth[scale=1.0]}, draw=blue!50!black},
    info/.style={-{Stealth[scale=1.0]}, draw=blue!50!black, dashed}
]
    % --- Node positions ---
    \node[chance]   (HD)   at ( 0,    5)    {HEARTDISEASE};
    \node[chance]   (ANG)  at (-4.5,  3.5)  {ANGIOGRAM};
    \node[chance]   (PAIN) at ( 4.5,  3.5)  {PAIN};
    \node[decision] (HS)   at (-4.5,  1.5)  {HEARTSURGERY};
    \node[chance]   (ER)   at ( 0,    1.5)  {EARLYRESULTS};
    \node[decision] (HP)   at ( 4.5,  1.5)  {HEARTPHARMA};
    \node[chance]   (LQ)   at (-3,   -0.5)  {LIFEQ};
    \node[chance]   (EC)   at ( 3,   -0.5)  {ECONOMICALC};
    \node[utility]  (UT)   at ( 0,   -2.6)  {UTILITY};

    % --- Probabilistic / functional arcs (solid) ---
    \draw[arr] (HD) -- (ANG);
    \draw[arr] (HD) -- (PAIN);
    \draw[arr] (HD) -- (LQ.north);
    \draw[arr] (ANG) -- (ER);
    \draw[arr] (HS) -- (ER);
    \draw[arr] (HS) -- (LQ);
    \draw[arr] (HS) -- (EC.west);
    \draw[arr] (HP) -- (EC);
    \draw[arr] (ER) -- (EC);
    \draw[arr] (LQ) -- (UT);
    \draw[arr] (EC) -- (UT);

    % --- Informational arcs (dashed, into decisions) ---
    \draw[info] (ANG) -- (HS);
    \draw[info] (ANG.south east) -- (HP.north west);
    \draw[info] (PAIN.south west) -- (HS.north east);
    \draw[info] (PAIN) -- (HP);
    \draw[info] (ER) -- (HP);
\end{tikzpicture}
\caption{Influence diagram of the \emph{bypass} model. Chance nodes
are ellipses, decision nodes rectangles and the utility node a
hexagon. Solid arcs are probabilistic or functional, dashed arcs
are informational.}
\label{fig:bypass-id}
\end{figure}

The Influence Diagram illustrated in Figure~\ref{fig:bypass-id} models a problem with two sequential decisions: first, \texttt{Heart Surgery (HS)}, and second, \texttt{Heart Pharma (HP)}. 

\texttt{Angiogram (ANG)} is a diagnostic test and \texttt{Pain (PA)} is a symptom; both variables are observed prior to making the decisions. \texttt{Early Results (ER)} is conditioned by \texttt{ANG} and \texttt{HS}, and constitutes known information for the \texttt{HP} decision. 

The preferences depend on \texttt{Life Quality (LQ)} and \texttt{Economic Calculation (EC)}. These are conditioned, respectively, by \texttt{HS} and \texttt{HD} (\texttt{Heart Disease)} for \texttt{LQ}; and by \texttt{ER}, \texttt{HS}, and \texttt{HP} for \texttt{EC}.

The model is regular. Variables \texttt{LQ} and \texttt{EC} have 3 levels, while the rest of the variables are binary. 

Regarding the parameterization, there are 9 utility parameters and $2 + 4 + 4 + 8 + 12 + 24 = 56$ parameters modeling the probability distributions. Consequently, there are $56 - 21 = 35$ degrees of freedom. In total, the model is defined by $35 + 9 = 44$ parameters.

\subsection{Quantitative parameterisation}

Each chance node $C_i$ carries a \emph{conditional probability
table} (CPT) $\theta_{C_i}$ that assigns to every configuration of
its probabilistic parents $\mathrm{Pa}(C_i)$ a probability distribution over
$\Omega_{C_i}$:
\[
   \theta_{C_i}\bigl(\cdot \mid \mathbf{pa}\bigr) \;\in\;
   \Delta(\Omega_{C_i})
   \qquad
   \forall \mathbf{pa}\in\prod_{C_j\in\mathrm{Pa}(C_i)}\Omega_{C_j},
\]
where $\Delta(\Omega_{C_i})$ denotes the probability simplex over the
states of $C_i$. The utility node carries a real-valued mapping
\[
   u : \prod_{X\in\mathrm{Pa}(V)}\Omega_X \longrightarrow \mathbb{R},
\]
where $\Omega_{D_i}\equiv\mathcal{A}_{D_i}$ for any decision parent.
The function $u$ encodes the preferences of the decision-maker in
the sense of Multi-Attribute Utility Theory, and when its arguments
satisfy the usual independence conditions it admits an additive or
multiplicative decomposition that drastically simplifies its
elicitation~\cite{riosinsua2002maut}. The collection of all CPTs
and the utility table is the \emph{parameter vector} of the
diagram, written $\boldsymbol{\Theta}$.

\begin{remark}[Combinatorial explosion]
\label{rem:explosion}
The number of free entries of a CPT grows multiplicatively with the
parents: a node with $m$ parents, each taking $r$ states, and $r$
own states contributes $r^{m}(r-1)$ free parameters. This
exponential growth is the source of the elicitation bottleneck
discussed in Chapter~\ref{ch:introduction} and the reason why
parameter optimisation rather than exhaustive elicitation is
necessary.
\end{remark}

A \emph{decision rule} for $D_i$ is a mapping
$\delta_i : \prod_{C_j\in\mathrm{Pa}(D_i)}\Omega_{C_j}
\longrightarrow \mathcal{A}_{D_i}$
that prescribes an action for every observable context. The
collection $\boldsymbol{\delta}=(\delta_1,\ldots,\delta_k)$ over all
decisions is called a \emph{policy}. The expected utility of a
policy under a parameter vector $\boldsymbol{\Theta}$ is
\begin{equation}
\label{eq:eu}
   \mathrm{EU}(\boldsymbol{\delta}\mid\boldsymbol{\Theta}) =
   \sum_{\mathbf{c}} u\!\bigl(\mathbf{c},
                              \boldsymbol{\delta}(\mathbf{c})\bigr)
   \prod_{C_i\in\mathbf{C}}
   \theta_{C_i}\!\bigl(c_i \mid \mathbf{pa}(c_i)\bigr),
\end{equation}
where $\mathbf{c}=(c_1,\ldots,c_n)$ is a complete \emph{scenario}. For each scenario,
$\mathbf{pa}(c_i)$ denotes the values its parents take within that same
scenario and $\boldsymbol{\delta}(\mathbf{c})$ the actions prescribed
by the policy. Thus each scenario is weighted by its joint probability
(the product of CPTs) and scored by its utility, and the expectation
sums these contributions over every possible state of the space.
The \emph{optimal policy} is
$\boldsymbol{\delta}^{\!*}(\boldsymbol{\Theta}) =
\arg\max_{\boldsymbol{\delta}}
\mathrm{EU}(\boldsymbol{\delta}\mid\boldsymbol{\Theta})$.

\noindent The time complexity of Shachter’s node-reduction and arc-reversal algorithm is actually exponential in the worst case, but 
the classical algorithm of
Shachter~\cite{shachter1986evaluating} computes
$\boldsymbol{\delta}^{\!*}(\boldsymbol{\Theta})$ in polynomial time relative to the size of the graph for any fixed $\boldsymbol{\Theta}$ by successively reversing arcs
and eliminating nodes from it.

\subsubsection*{What the parameters and their evaluation look like}
\label{sec:bypass-tables}

The three pieces of the parameterisation become clear when one
looks at the tables of the bypass model of
Figure~\ref{fig:bypass-id}. Table~\ref{tab:bypass-cpt} shows a real
CPT extracted from the model: $P(\texttt{Angiogram}\mid
\texttt{HeartDisease})$, which encodes the standard
sensitivity/specificity of the angiogram test.

\begin{table}[h]
\centering\small
\begin{tabular}{l|cc}
\hline
$P(\texttt{Angiogram} \mid \texttt{HeartDisease})$
              & \textsc{negative} & \textsc{positive}\\
\hline
$\texttt{HeartDisease} = \textsc{absent}$  & 0.95 & 0.05\\
$\texttt{HeartDisease} = \textsc{present}$ & 0.14 & 0.86\\
\hline
\end{tabular}
\caption{Real CPT extracted from the bypass model.
Each row is a probability distribution over the two test outcomes.}
\label{tab:bypass-cpt}
\end{table}

The utility node \texttt{Utility} of the bypass diagram depends on
the quality-of-life node \texttt{LifeQ} and the cost node
\texttt{Economicalc}; Table~\ref{tab:bypass-util} shows its full set
of values, taken directly from the model. As one would expect, the
best outcome (high quality of life, low cost) takes the maximum
value, while death together with high cost takes the minimum.

\begin{table}[h]
\centering\small
\begin{tabular}{l|ccc}
\hline
$u(\texttt{LifeQ}, \texttt{Economicalc})$
                                 & \textsc{low} & \textsc{medium} & \textsc{high}\\
\hline
$\texttt{LifeQ} = \textsc{dead}$     & 2.00 & 0.70 & 0.05\\
$\texttt{LifeQ} = \textsc{live2alq}$ & 3.10 & 2.90 & 2.80\\
$\texttt{LifeQ} = \textsc{live2ahq}$ & 7.80 & 4.80 & 2.80\\
\hline
\end{tabular}
\caption{Utility table of the bypass model.}
\label{tab:bypass-util}
\end{table}

Given the two tables above (and the rest of the CPTs of the model,
not shown), the Shachter algorithm produces a third one: the
expected utility of each action of a decision node \emph{conditioned
on the information available when that decision is taken}.
Table~\ref{tab:bypass-eu} shows this output for the
\texttt{HeartSurgery} decision, whose informational parents are
\texttt{Pain} and \texttt{Angiogram}.

\begin{table}[h]
\centering\small
\begin{tabular}{ll|cc}
\hline
\texttt{Pain} & \texttt{Angiogram}
              & $\mathrm{EU}(\textsc{no})$
              & $\mathrm{EU}(\textsc{yes})$\\
\hline
\textsc{absent}  & \textsc{negative} & \textbf{6.98} & 3.73 \\
\textsc{absent}  & \textsc{positive} & \textbf{5.34} & 3.91 \\
\textsc{present} & \textsc{negative} & \textbf{6.66} & 3.78 \\
\textsc{present} & \textsc{positive} & 3.43          & \textbf{4.18}\\
\hline
\end{tabular}
\caption{Expected utility (EU) of each action of \texttt{HeartSurgery}
under the parameters of the bypass model, as produced by the
Shachter evaluation. The bold value in each row is the optimal
action for that information state. Read as a decision rule:
\emph{operate only when pain is present and the angiogram is
positive; abstain otherwise}.}
\label{tab:bypass-eu}
\end{table}

Schematically, the standard use of an ID flows from CPTs and the
utility function towards the optimal policy:
\[
  (\theta_{C_i},\, u) \xrightarrow{\text{Shachter}} \mathrm{EU}
  \xrightarrow{\arg\max} \boldsymbol{\delta}^{\!*}.
\]
The problem this
thesis addresses runs the arrows backwards. The expert does
\emph{not} supply Tables~\ref{tab:bypass-cpt}
and~\ref{tab:bypass-util}; he supplies only a few entries of
Table~\ref{tab:bypass-eu} as \emph{decision rules} (for example:\textit{ do not operate
when pain is absent and the angiogram is negative}). The optimizer
must recover a pair $(\theta_{C_i},\, u)$ \emph{compatible} with
those rules, in the sense that re-running Shachter on the recovered
parameters reproduces the expert's choices.

Preference is represented using lotteries for comparisons between uncertain consequences, for example:
\[
(\texttt{LQ} = \text{Dead}, \text{EC} = \text{High}; 0.25) \leq (\texttt{LQ} = \text{Live2AHQ}, \texttt{EC} = \text{Low}; 7.5)
\]

The very possibility of summarizing the expert's preferences by a
single utility function is not gratuitous: it rests on the classical
axioms of rational choice.



\begin{definition}[von Neumann--Morgenstern axioms]
\label{def:vnm-axioms}
A preference relation $\succeq$ over lotteries (probability
distributions over outcomes) admits a representation by
\emph{expected utility} (i.e.\ there exists a function $u$ such that
$L_1 \succeq L_2 \iff \mathbb{E}_{L_1}[u] \ge \mathbb{E}_{L_2}[u]$)
\emph{if and only if} it satisfies the following
axioms~\cite{vonneumann1944theory}:
\begin{itemize}
  \item \textbf{Completeness (orderability).} For any two lotteries
        $L_1$ and $L_2$, exactly one of $L_1 \succ L_2$,
        $L_2 \succ L_1$ or $L_1 \sim L_2$ holds; the agent is never
        unable to compare two alternatives.
  \item \textbf{Transitivity.} If $L_1 \succeq L_2$ and
        $L_2 \succeq L_3$, then $L_1 \succeq L_3$.
  \item \textbf{Continuity.} If $L_1 \succeq L_2 \succeq L_3$, there
        exists a probability $p \in [0,1]$ such that the agent is
        indifferent between $L_2$ and the mixed lottery
        $p\,L_1 + (1-p)\,L_3$.
  \item \textbf{Independence (substitutability).} If
        $L_1 \succeq L_2$, then for any third lottery $L_3$ and any
        $p \in (0,1]$ we have
        $p\,L_1 + (1-p)\,L_3 \succeq p\,L_2 + (1-p)\,L_3$; mixing both
        options with a common third lottery cannot reverse the
        preference.
\end{itemize}
\end{definition}

The relevance of the axioms here is conceptual: whenever the
recovered pair $(\theta_{C_i},\, u)$ induces a policy through the
maximum-expected-utility rule, that policy is \emph{rational} in the
von Neumann--Morgenstern sense by construction. The expert rules, by
contrast, carry no such guarantee: they are isolated \emph{if--then}
prescriptions that need not derive from any coherent underlying
preference order.

This gap raises what we regard as a central theoretical open question
behind the whole recovery program:

\begin{quote}\itshape
It remains to be seen whether the algorithm can find, for \emph{any} set of expert rules, a parameterization $(\theta_{C_i},\, u)$ that is simultaneously consistent with those
rules and with the axioms of utility theory; or whether some rule
sets are intrinsically incoherent, admitting no rational Influence
Diagram able to reproduce them.
\end{quote}

% ==============================================================
\section{The Knowledge-Acquisition Problem}
% ==============================================================

In practice the structure of the ID is elicited from the expert in a
small number of modelling sessions, but the numerical
parameterisation $\boldsymbol{\Theta}$ is far harder to obtain. The
classical cycle of~\cite{bielza2000issues,bielza2010modeling} proceeds in
five stages — problem identification, enumeration of alternatives,
modelling, evaluation and sensitivity analysis — and the modelling
stage alone accounts for the bulk of the cost. As anticipated in
Remark~\ref{rem:explosion}, the size of $\boldsymbol{\Theta}$ grows
exponentially with the connectivity of the diagram, so an exhaustive
elicitation is unrealistic for any diagram of practical size.

We assume instead that the expert can provide a \emph{small set of
decision rules}.

\begin{definition}[Expert rule base]
\label{def:rules}
A \emph{rule base} is a finite set
\[
   \mathcal{R} = \bigl\{(\mathbf{x}^{(l)}, D_{i_l}, a^{(l)})\bigr\}_{l=1}^{L},
\]
where $\mathbf{x}^{(l)}$ is a (possibly partial) observation vector
on the informational parents of decision $D_{i_l}$, and
$a^{(l)}\in\mathcal{A}_{D_{i_l}}$ is the action the expert deems
optimal under $\mathbf{x}^{(l)}$. Components marked as
``irrelevant'' contribute no constraint and the rule must hold for
every value they may take.
\end{definition}

Given such a rule base we measure the consistency of a parameter
vector with the expert by a simple accuracy:

\begin{definition}[Fitness function]
\label{def:fitness}
Define the vector of expected utilities of decision $D_{i_l}$ under
$\boldsymbol{\Theta}$ and evidence $\mathbf{x}^{(l)}$ as
\[
   \mathbf{u}^{(l)}(\boldsymbol{\Theta}) =
   \bigl(\mathrm{EU}(a \mid \boldsymbol{\Theta}, \mathbf{x}^{(l)})\bigr)_{a\in\mathcal{A}_{D_{i_l}}}.
\]
The generic loss minimised by the EDA is then
\[
   \mathcal{L}(\boldsymbol{\Theta}) =
   \sum_{(\mathbf{x}^{(l)},\,D_{i_l},\,a^{(l)})\in\mathcal{R}}
   \ell\!\left(\mathbf{u}^{(l)}(\boldsymbol{\Theta}),\; a^{(l)}\right),
\]
where $\ell$ is a per-rule loss that compares the expected-utility
vector of the target decision with the action prescribed by the
expert. Different choices of $\ell$ recover the fitness variants used
in this work:
\begin{itemize}
   \item \textbf{Binary} ($0$--$1$ loss):
         $\ell_{\mathrm{binary}}(\mathbf{u}, a) = \mathbf{1}\!\left[\textstyle\arg\max_{a'} u_{a'} \neq a\right]$.
   \item \textbf{Regret}:
         $\ell_{\mathrm{regret}}(\mathbf{u}, a) = \max_{a'} \left( u_{a'} - u_a \right)$.
   \item \textbf{Margin}:
         $\ell_{\mathrm{margin}}(\mathbf{u}, a) = \max\!\bigl(0,\; \max_{a'\neq a} u_{a'} + m - u_a\bigr)$.
   \item \textbf{Softmax} (cross-entropy):
         $\ell_{\mathrm{softmax}}(\mathbf{u}, a) = -\log \operatorname{softmax}(\mathbf{u})_a$.
   \item \textbf{Entropy-regularized}:
         $\ell_{\mathrm{entropy}}(\mathbf{u}, a) = -\log \operatorname{softmax}(\mathbf{u})_a + \alpha\, H\!\bigl(\operatorname{softmax}(\mathbf{u})\bigr)$.
\end{itemize}
where $m$ is the margin, $\alpha$ the entropy-regularization weight, and
$H(\mathbf{p}) = -\sum_{a'} p_{a'}\log p_{a'}$ the Shannon entropy of the
softmax distribution over actions.

\end{definition}

The parameter-recovery problem is then the constrained search
\[
   \boldsymbol{\Theta}^{\star} =
   \arg\min_{\boldsymbol{\Theta}} \mathcal{L}(\boldsymbol{\Theta}),
\]
with $\boldsymbol{\Theta}$ ranging over the product of all valid
CPTs and utility tables. The next proposition justifies the use of
a derivative-free metaheuristic instead of a gradient-based
optimiser.

\begin{proposition}[Irregularity of the objective]
\label{prop:nonsmooth}
The objective $\mathcal{L}$ of Definition~\ref{def:fitness} is
non-convex with respect to $\boldsymbol{\Theta}$ for every per-rule
loss $\ell$ considered in this work. For the $0$--$1$ (\textsc{binary})
loss it is, in addition, piecewise constant and discontinuous; for the
\textsc{regret} and \textsc{margin} losses it is continuous but
non-differentiable on the set where the optimal action of a decision
switches; for \textsc{softmax} and \textsc{entropy} it is piecewise
smooth.
\end{proposition}

\begin{proof}[Sketch]
\textbf{\emph{Non-convexity.}} For a chain of two chance nodes the expected
utility contains a term proportional to the product of two parameters,
$p\,q$. The map $(p,q)\mapsto p\,q$ is not convex: at $(1,0)$ and
$(0,1)$ it equals $0$, yet at the midpoint $(\tfrac12,\tfrac12)$ it
equals $\tfrac14>0$, violating the convexity inequality. Since
$\mathbf{u}^{(l)}(\boldsymbol{\Theta})$ generically involves such
products (a node multiplies the probabilities of its parents) and a
normalisation by the probability of the evidence, $\mathcal{L}$ is
non-convex.

\textbf{\emph{Regularity.}} Within any region of $\boldsymbol{\Theta}$ where the
optimal action of every decision is fixed, $\mathbf{u}^{(l)}$ is a
rational function of the parameters and is therefore smooth. At the
boundaries between regions the maximising action switches: the value of the optimal policy is non-differentiable there. Composing with the per-rule loss,
the $0$--$1$ loss jumps discontinuously at those boundaries, the
\textsc{regret} and \textsc{margin} losses stay continuous but
non-differentiable, and the \textsc{softmax} and \textsc{entropy} losses
are smooth away from them.
\end{proof}

\noindent A gradient of $\mathcal{L}$ thus exists almost everywhere
(except for the discontinuous \textsc{binary} loss) and could in
principle be obtained by finite differences. We nevertheless adopt a
population-based, derivative-free strategy for three reasons: 

\begin{itemize}
    \item The landscape is strongly non-convex and highly degenerate (many parameter vectors reproduce the same rules) so local gradient information does not
lead to the global optimum.
    \item The evaluation engine returns only utility
values, so each numerical gradient would cost $O(d)$ additional
evaluations in a space of dimension $d=|\boldsymbol{\Theta}|$.
    \item And for the
\textsc{binary} loss the gradient is identically zero almost everywhere.
\end{itemize}


% ==============================================================
\section{Estimation of Distribution Algorithms}
% ==============================================================

Estimation of Distribution Algorithms
(EDAs)~\cite{larranaga2002eda,larranaga2024edas} are a class of
evolutionary metaheuristics in which the classical
crossover-and-mutation operators are replaced by the
\emph{construction and sampling of an explicit probabilistic model}
of the search space. At every generation a probability distribution
is fitted to the best individuals of the current population and used
to generate the next generation. Algorithm~\ref{alg:eda} states the
generic scheme.

\begin{algorithm}[Generic EDA]\hfill
\label{alg:eda}
\begin{algorithmic}[1]
\State Initialise a population $\mathcal{P}_0$ of $N$ individuals from a prior on the search space.
\State $t \leftarrow 0$
\Repeat
    \State Evaluate $f$ on every $\boldsymbol{\Theta}\in\mathcal{P}_t$.
    \State $\mathcal{P}_t^{\text{sel}} \leftarrow$ top-$S$ individuals of $\mathcal{P}_t$.
    \State Estimate a probabilistic model $\hat{p}_t$ based on
           $\mathcal{P}_t^{\text{sel}}$.
    \State Sample $N$ new individuals from $\hat{p}_t$ to form
           $\mathcal{P}_{t+1}$.
    \State $t \leftarrow t+1$
\Until{stopping criterion is met}
\State \Return best individual found.
\end{algorithmic}
\end{algorithm}

In words, the algorithm keeps a \emph{population} of candidate solutions and
repeatedly (i)~scores each one with the fitness $f$, (ii)~keeps the top-$S$
elite, (iii)~fits a probabilistic model $\hat{p}_t$ to that elite, and
(iv)~samples the next generation from it. It thus replaces the crossover and
mutation of a genetic algorithm by \emph{learning and sampling a distribution}
over the good solutions, which is what lets it exploit statistical regularities
of the search space.

A candidate solution (an \emph{individual}) is a real-valued vector
$\boldsymbol{\phi}=(\phi_1,\ldots,\phi_d)\in\mathbb{R}^d$ whose components stack,
in a fixed order, the raw value of every free CPT entry and every utility entry
of the diagram. The search itself is \emph{unconstrained} — each $\phi_j$ lives in the whole
real line $\mathbb{R}$, and the bounds $[-5,5]$/$[-10,10]$ constrain only the
initial sampling — and it is a decoder
(Chapter~\ref{ch:implementation}) that turns $\boldsymbol{\phi}$ into a valid
parameter vector $\boldsymbol{\Theta}$, enforcing the two constraints of the
model: every CPT row must be a probability distribution (non-negative, summing to
one) and every utility must lie inside the expert range $[u_{\min},u_{\max}]$.
As an example, for a diagram with a single binary chance node conditioned on a
binary parent ($2$ rows of $2$ entries) and $2$ utility values, one valid
individual is
\[
  \boldsymbol{\phi} =
  \bigl(\underbrace{1.2,\,-0.4}_{\text{CPT row }1},\;
        \underbrace{0.1,\,0.3}_{\text{CPT row }2},\;
        \underbrace{2.0,\,-1.5}_{\text{utilities}}\bigr),
\]
which the decoder maps (row-wise softmax and a rescaled sigmoid) to, e.g.,
CPT rows $(0.83,0.17)$ and $(0.45,0.55)$ and utilities $(8.8,\,1.8)$ on the range
$[0,10]$ — a fully valid model regardless of the raw values sampled.

The defining choice of an EDA is the family used for $\hat{p}_t$:
richer families capture more structure of the search space at the
expense of a higher estimation cost.

\subsection{Dependency Models in EDAs}

Because the parameter vector of an ID lives in a continuous space
(every CPT entry and every utility value is a real number), only
continuous EDAs are relevant. The four models used in this work,
all available through \textsc{EDAspy}~\cite{soloviev2024edaspy}, span the spectrum from a fully factorised parametric model to a
full-covariance multivariate Gaussian, with a structured intermediate
and a non-parametric alternative on the univariate side. We denote an individual by $\boldsymbol{\phi}=(\phi_1,\ldots,\phi_d)$
and the size of the selected elite by
$S=|\mathcal{P}_t^{\text{sel}}|$.

\paragraph*{UMDA\textsubscript{c} --- Univariate Marginal Distribution
Algorithm.} Each dimension is modeled by an independent
Gaussian~\cite{larranaga2002eda}:
\[
    \hat{p}_t(\boldsymbol{\phi})
    = \prod_{j=1}^{d}
    \mathcal{N}\!\bigl(\phi_j;\,
    \hat{\mu}_t^{j},\,(\hat{\sigma}_t^{j})^{2}\bigr),
\]
whose mean and variance are the sample moments of the selected individuals,
\[
    \hat{\mu}_t^{j}
    = \frac{1}{S}\!\sum_{\boldsymbol{\phi}\in\mathcal{P}_t^{\text{sel}}}\!\phi_j,
    \qquad
    (\hat{\sigma}_t^{j})^{2}
    = \frac{1}{S}\!\sum_{\boldsymbol{\phi}\in\mathcal{P}_t^{\text{sel}}}\!
    (\phi_j - \hat{\mu}_t^{j})^{2}.
\]
The model is the cheapest of the family ($O(d\cdot S)$ per generation,
both for fitting and for sampling) but its independence assumption
means it cannot capture any dependency between the parameters of the
diagram: the search effectively assumes that good values for one CPT
entry give no information about good values for another.
\paragraph*{EMNA --- Estimation of Multivariate Normal Algorithm.}
The opposite end of the parametric spectrum: the whole population is
modelled by a single multivariate Gaussian~\cite{larranaga2002eda},
\[
    \hat{p}_t(\boldsymbol{\phi}) =
    \mathcal{N}\!\bigl(\boldsymbol{\phi};\,
    \hat{\boldsymbol{\mu}}_t,\,
    \hat{\boldsymbol{\Sigma}}_t\bigr),
\]
with sample mean and \emph{full covariance},
\[
    \hat{\boldsymbol{\mu}}_t
    = \frac{1}{S}\!\sum_{\boldsymbol{\phi}\in\mathcal{P}_t^{\text{sel}}}\!\boldsymbol{\phi},
    \qquad
    \hat{\boldsymbol{\Sigma}}_t
    = \frac{1}{S}\!\sum_{\boldsymbol{\phi}\in\mathcal{P}_t^{\text{sel}}}\!
    (\boldsymbol{\phi} - \hat{\boldsymbol{\mu}}_t)
    (\boldsymbol{\phi} - \hat{\boldsymbol{\mu}}_t)^{\!\top}.
\]
EMNA captures every pairwise linear dependency between parameters and
is therefore strictly more expressive than UMDA. The price is
double: the cost grows to $O(d^{2}\cdot S)$ for the covariance estimate,
plus a $O(d^{3})$ Cholesky factorisation of $\hat{\Sigma}_t$ used to sample (and the
truncated population must be large enough, $S>d$
to keep $\hat{\boldsymbol{\Sigma}}_t$
non-singular)~\cite{larranaga2002eda}.

\paragraph*{EGNA --- Estimation of Gaussian Networks Algorithm.}
A middle ground between UMDA and EMNA: a multivariate Gaussian whose
factorisation is governed by a Bayesian network $\mathcal{G}$
\emph{learnt from the elite at every
generation}~\cite{larranaga2002eda}. Each variable conditions only
on its parents in $\mathcal{G}$,
\[
    \hat{p}_t(\boldsymbol{\phi}) = \prod_{j=1}^{d}
    \mathcal{N}\!\Bigl(\phi_j;\;
    \beta_{j,0} + \!\!\!\!\sum_{k\,\in\,\mathrm{pa}_{\mathcal{G}}(j)}\!\!\!\!
    \beta_{j,k}\,\phi_k,\;
    \sigma_j^{2}\Bigr),
\]
i.e.\ a linear regression of $\phi_j$ on its graph parents with
Gaussian residuals. The structure $\mathcal{G}$ is chosen with a
score (typically BIC) and a greedy search over
DAGs~\cite{larranaga2002eda}. EGNA recovers UMDA when $\mathcal{G}$
has no edges and approaches EMNA when $\mathcal{G}$ becomes a
complete DAG; in practice the learnt graph is sparse, which keeps
the cost reduced and makes the model robust if 
$S$ is not large enough for EMNA to estimate a full covariance
reliably.

\paragraph*{KEDA (Univariate) --- Kernel-density Estimation of
Distribution Algorithm.} A non-parametric alternative on the
univariate side: each marginal is fitted with a kernel density
estimator from the elite~\cite{soloviev2024edaspy},
\[
    \hat{p}_t(\boldsymbol{\phi}) = \prod_{j=1}^{d}
    \hat{p}_t^{\,j}(\phi_j),
    \qquad
    \hat{p}_t^{\,j}(\phi_j)
    = \frac{1}{S\,h_j}\!\sum_{\boldsymbol{\phi}^{(s)}\in\mathcal{P}_t^{\text{sel}}}\!\!
    K\!\Bigl(\frac{\phi_j - \phi_j^{(s)}}{h_j}\Bigr),
\]
where $K$ is a kernel (typically Gaussian) and $h_j$ is a bandwidth
chosen, for instance, by Silverman's rule. Like UMDA it ignores
cross-parameter dependencies, but it drops the Gaussian assumption
per dimension and can therefore fit skewed or multi-modal marginals
(a setting in which a single Gaussian would smear over distinct
modes). Its cost is again $O(d\cdot S)$, with the bandwidth as an
additional implicit hyperparameter that trades smoothness against
detail.

\medskip

Together the four variants span the spectrum from \emph{fully
factored} (UMDA\textsubscript{c}, KEDA) to \emph{fully joint} (EMNA),
with EGNA as a structured intermediate, and from \emph{parametric}
(UMDA\textsubscript{c}, EMNA, EGNA) to \emph{non-parametric} (KEDA)
on the per-dimension side. Which combination suits the
parameter-recovery problem of this work is far from obvious a
priori.

\subsection{Convergence of the Algorithm}

The reason an EDA makes progress is intuitive and needs no heavy
machinery. At each generation it keeps only the best fraction of the
population — a \emph{truncation} selection that retains the top
$\mu\in(0,1)$ of the individuals — and refits the probabilistic model
$\hat{p}_t$ to that elite. Because the selected individuals
concentrate in the better-scoring regions of the search space, the
model fitted to them places its mass there, so the next population,
sampled from $\hat{p}_t$, is on average better than the previous one.

Repeating the loop has two coupled effects. The model \emph{drifts}
towards regions of lower cost, generation after generation, while its
variance \emph{contracts} as the surviving individuals become more
alike. The drift is what improves the solutions; the contraction is
what makes the search eventually settle on a concrete configuration.
When the variance has shrunk to the point where new samples barely
differ from one another, the population has converged.

The selection ratio $\mu$ balances these two effects. A small $\mu$
keeps only a handful of top individuals: the model contracts quickly
and the search converges fast, but it may lock onto the first good
region it finds. A large $\mu$ keeps almost the whole population, so
little information is extracted from the cost function and the search
behaves almost like random sampling. A useful value lies in between
and is tuned empirically in Chapter~\ref{ch:experiments}.

This argument is deliberately informal. It explains why the
distribution concentrates on good regions, not that the region found
is the global optimum — on the non-convex, degenerate landscape of
Proposition~\ref{prop:nonsmooth} it generally need not be. What
matters in practice is that the EDA reliably drives the cost down
towards high-quality solutions, which is what the experiments of
Chapter~\ref{ch:experiments} confirm. The only additional requirement specific to our
problem is that every sampled individual be turned into a valid set of
CPTs — each row a probability distribution — which is handled at the
decoding stage described in Chapter~\ref{ch:implementation}.

The solution reached must be validated by the expert and consider several iterations, revising/expanding the input rule base, adding constraints to the parameter space, and possibly revising the dependency graph.